\chapter{微分形式与斯托克斯定理}\label{ch:b3:forms}

两个世纪以来，分析学中的一条定理
吸收了同类的其他全部定理：微积分基本定理、
格林–黎曼公式（证明见本科二年级分册）、高斯散度定理、
开尔文–斯托克斯旋度定理——它们每一条都断言，
某个导数在区域上的积分，
等于原对象在边界上的积分。以\emph{微分形式}为语言，
把它们化为同一条命题，
$\int_M\dd\omega = \int_{\partial M}\omega$，并使该命题
可一气呵成地证明。本章将
严谨地建立这一语言：交错多重线性代数、
外微分、拉回、定向，以及
在\cref{ch:b3:submanifolds}中的子流形上的积分——
证明斯托克斯定理，并兑现早先的承诺：经典积分
定理，以及暗中支撑了
\cref{ch:b3:residues}中的环绕数，以及周末问题中的布劳威尔
不动点定理。下文中，\emph{光滑}表示 $\mathcal
C^\infty$；除非另有说明，所有映射与形式都是光滑的。
这在本阶段不损失值得保留的一般性，
却使论证大为简便。

\section{交错多重线性代数}

\begin{definition}\label{def:b3:forms:alternating}
设 $E$ 为 $n$ 维实向量空间。称
\emph{$k$ 重线性交错形式}\index{交错形式}为
$E$ 上的映射 $\alpha\colon E^k \to \R$，关于每个
变元线性，且只要两个变元相等，$\alpha(v_1, \dots, v_k) = 0$。
这些形式所成的空间记为 $\Lambda^k E^*$；
约定 $\Lambda^0E^* = \R$。交错性蕴含
反对称性：交换两个变元会改变符号
（展开 $\alpha(\dots, v + w, \dots, v + w, \dots) = 0$），
并且一般地，$\alpha(v_{\sigma(1)}, \dots, v_{\sigma(k)}) =
\varepsilon(\sigma)\,\alpha(v_1, \dots, v_k)$，对每个
置换 $\sigma$ 成立。
\end{definition}

\begin{example}
在 $E = \R^n$ 上，$\Lambda^1E^* = E^*$ 即对偶空间；
标准基下的行列式是一个 $n$ 重线性
交错形式，\cref{prop:b3:forms:basis}将说明它
张成 $\Lambda^nE^*$——这正是行列式
在相差一个非零倍数之意义下唯一的深层原因。当 $k > n$ 时，$\Lambda^kE^* = \{0\}$：
$k$ 个向量线性相关，将其中一个沿其余向量展开，
由交错性即得 $\alpha$ 为零。
\end{example}

\begin{definition}\label{def:b3:forms:wedge}
对 $\ell_1, \dots, \ell_k \in E^*$，称它们的\emph{外
积}\index{外积}为如下的 $k$ 重线性交错
形式：
\[
(\ell_1 \wedge \dots \wedge \ell_k)(v_1, \dots, v_k)
= \det\bigl(\ell_i(v_j)\bigr)_{1 \leq i, j \leq k} .
\]
交错性与多重线性都继承自行列式按
其列所具有的相应性质。
\end{definition}

\begin{proposition}[$\Lambda^k$ 的基]\label{prop:b3:forms:basis}
设 $(e_1, \dots, e_n)$ 为 $E$ 的一组基，其对偶基为
$(e_1^*, \dots, e_n^*)$。下述形式
\[
e_I^* = e_{i_1}^* \wedge \dots \wedge e_{i_k}^*,
\qquad I = \{i_1 < \dots < i_k\} \subseteq \{1, \dots, n\},
\]
构成 $\Lambda^kE^*$ 的一组基；由此 $\dim\Lambda^kE^* =
\binom nk$。具体地，$\alpha =
\sum_{\abs I = k}\alpha(e_{i_1}, \dots, e_{i_k})\,e_I^*$。
\end{proposition}

\begin{proof}
\emph{生成性。}设 $\alpha \in \Lambda^kE^*$，且 $\beta =
\sum_I\alpha(e_I)\,e_I^*$，其中 $\alpha(e_I)$ 是
$\alpha(e_{i_1}, \dots, e_{i_k})$ 的缩写。两边都是 $k$ 重线性
且交错，因此只要两者在所有
$k$ 元组 $(e_{j_1}, \dots, e_{j_k})$，其中 $j_1 < \dots <
j_k$（多重线性把问题归结为基向量组，
交错性则归结为严格递增的情形）。而 $e_I^*(e_{j_1},
\dots, e_{j_k}) = \det(e_{i_r}^*(e_{j_s})) = \delta_{IJ}$：
当 $I = J$ 时矩阵为单位阵；当 $I \neq J$ 时某一行
为零。于是 $\beta(e_J) = \alpha(e_J)$ 对一切 $J$ 成立：
$\beta = \alpha$。\emph{线性无关性。}若 $\sum_I c_Ie_I^* = 0$，
在 $(e_{j_1}, \dots, e_{j_k})$ 上求值即得 $c_J = 0$。
\end{proof}

\begin{definition}\label{def:b3:forms:wedgegeneral}
外积可延拓为一个双线性映射
$\Lambda^kE^* \times \Lambda^\ell E^* \to
\Lambda^{k+\ell}E^*$，由双线性性以及
$(e_I^*) \wedge (e_J^*) = e_I^* \wedge e_J^*$（拼接后重排；
且乘积为 $0$，当 $I \cap J \neq
\varnothing$ 时）。该运算满足结合律，
并具有\emph{分次反交换性}：
\[
\beta \wedge \alpha = (-1)^{k\ell}\,\alpha \wedge \beta
\qquad (\alpha \in \Lambda^k, \ \beta \in \Lambda^\ell).
\]
\end{definition}

\begin{proof}[\omnameProof]
两条性质都先在基元素上验证，再由
双线性性推广。关于结合律，
$e_I^* \wedge e_J^* \wedge e_K^*$ 的两种加括号方式都等于由
$1$-形式拼接而成的楔积，依据
\cref{def:b3:forms:wedge}中的行列式公式（按块作拉普拉斯展开）。
符号规则：把 $\ell$ 个 $\beta$ 因子中的每一个
越过 $k$ 个 $\alpha$ 因子，每次相邻对换产生一个符号，
（即行列式两行的互换），因此总共
得到 $(-1)^{k\ell}$。
\end{proof}

\begin{proposition}[拉回，线性情形]
\label{prop:b3:forms:linearpullback}
线性映射 $u\colon E \to F$ 对每个 $k$ 诱导出
线性映射 $u^*\colon \Lambda^kF^* \to \Lambda^kE^*$，
$(u^*\alpha)(v_1, \dots, v_k) = \alpha(u(v_1), \dots,
u(v_k))$。它满足 $u^*(\alpha \wedge \beta) = u^*\alpha
\wedge u^*\beta$，且 $(u \circ w)^* = w^* \circ u^*$。
此外：
\begin{enumerate}
  \item 若 $\dim E = n$ 且 $u\colon E \to E$，则在
        直线 $\Lambda^nE^*$ 上：$u^*\alpha = (\det u)\,\alpha$。
  \item 若 $\operatorname{rk}u < k$，则 $u^* = 0$ 在
        $\Lambda^kF^*$ 上。
\end{enumerate}
\end{proposition}

\begin{proof}
函子性恒等式可直接由定义得出
（对乘积法则，在 $1$-形式的楔积上
用行列式公式验证 --- $\det(\ell_i(uv_j)) =
\det((u^*\ell_i)(v_j))$ --- 再双线性扩张）。
(1) $u^*$ 把一维的 $\Lambda^nE^*$
（\cref{prop:b3:forms:basis}）映到自身，故 $u^*\alpha =
c\,\alpha$，其中 $c$ 与 $\alpha \neq 0$ 无关；取
$\alpha = e_1^* \wedge \dots \wedge e_n^*$ 及 $(v_j) =
(e_j)$ 检验，得 $c = \det(e_i^*(ue_j)) = \det u$。
(2) 对 $v_1, \dots, v_k \in E$，向量 $u(v_1), \dots,
u(v_k)$ 位于 $u$ 的像中，维数 $< k$：它们
线性相关，交错形式在相关族上
取值为零（把相关向量沿
其余向量展开）。
\end{proof}

\section{微分形式与外微分}

\begin{definition}\label{def:b3:forms:diffform}
设 $U \subseteq \R^n$ 为开集。称 \emph{微分
$k$-形式}\index{微分形式}为 $U$ 上的光滑映射
$\omega\colon U \to \Lambda^k(\R^n)^*$；在
\cref{prop:b3:forms:basis} 的基下（记 $\dd x_i$ 为
$e_i^*$），
\[
\omega = \sum_{\abs I = k} a_I\,\dd x_I,
\qquad
\dd x_I = \dd x_{i_1} \wedge \dots \wedge \dd x_{i_k},
\]
其中系数 $a_I \in \mathcal C^\infty(U)$ 均为光滑函数。
其全体记为 $\Omega^k(U)$；$\Omega^0(U) = \mathcal
C^\infty(U)$。$0$-形式即函数；$1$-形式是
线性形式的场（例如函数的微分 $\dd f$
）；$n$-形式为 $a\,\dd x_1\wedge\dots\wedge\dd
x_n$，即 \cref{ch:b3:product} 中的自然被积对象。
\end{definition}

\begin{definition}[外微分]\label{def:b3:forms:d}
\emph{外微分}\index{外微分}是
线性映射 $\dd\colon \Omega^k(U) \to \Omega^{k+1}(U)$，
由下式定义：
\[
\dd\Bigl(\sum_I a_I\,\dd x_I\Bigr)
= \sum_I \dd a_I \wedge \dd x_I
= \sum_I\sum_{j=1}^n \frac{\partial a_I}{\partial
x_j}\,\dd x_j \wedge \dd x_I .
\]
对 $0$-形式，它即通常的微分。
\end{definition}

\begin{theorem}\label{thm:b3:forms:dsquare}
(a) $\dd(\omega \wedge \eta) = \dd\omega \wedge \eta +
(-1)^k\,\omega \wedge \dd\eta$，对 $\omega \in \Omega^k$
（即 \emph{分次莱布尼茨法则}）。
(b) $\dd \circ \dd = 0$。\index{$\dd^2 = 0$}
\end{theorem}

\begin{proof}
(a) 由双线性性，只需考虑 $\omega = a\,\dd
x_I$，$\eta = b\,\dd x_J$。于是 $\omega \wedge \eta =
ab\,\dd x_I \wedge \dd x_J$，且
\[
\dd(\omega\wedge\eta) = (b\,\dd a + a\,\dd b) \wedge \dd x_I
\wedge \dd x_J
= (\dd a \wedge \dd x_I) \wedge (b\,\dd x_J) +
a\,\dd b \wedge \dd x_I \wedge \dd x_J,
\]
并将 $1$-形式 $\dd b$ 越过 $k$ 个 $\dd
x_I$ 的因子，产生符号 $(-1)^k$（\cref{def:b3:forms:wedgegeneral}）：
第二项为 $(-1)^k\,\omega \wedge \dd\eta$。
(b) 对函数，$\dd(\dd a) = \sum_{i,j}
\frac{\partial^2a}{\partial x_j\partial x_i}\,\dd x_j \wedge
\dd x_i$ 的系数关于 $(i,j)$ 对称，由
施瓦茨定理（关于混合偏导数，在本科二年级分册中证明，
$a$ 为 $\mathcal C^\infty$），而 $\dd x_j \wedge
\dd x_i$ 反对称：把 $(i,j)$ 项与
$(j,i)$ 项配对，全部相消。对一般的 $\omega = \sum
a_I\dd x_I$：$\dd\dd\omega = \sum\dd(\dd a_I \wedge \dd x_I)
= \sum(\dd\dd a_I\wedge\dd x_I - \dd a_I\wedge\dd(\dd x_I))$
由 (a)，两项均为零（$\dd(\dd x_I) = 0$，因为
系数为常数）。
\end{proof}

\begin{definition}[拉回]\label{def:b3:forms:pullback}
设 $\varphi\colon U \to V$ 为光滑映射（$U \subseteq \R^m$，
$V \subseteq \R^n$ 为开集）。称 \emph{拉回}\index{拉回}
$\varphi^*\colon \Omega^k(V) \to \Omega^k(U)$ 为映射，它
逐点由沿微分的线性拉回定义：
$(\varphi^*\omega)_x = (D\varphi(x))^*\,\omega_{\varphi(x)}$。
具体地，$\varphi^*$ 作代换：对函数，$\varphi^*f = f \circ
\varphi$；$\varphi^*(\dd y_i) =
\dd\varphi_i = \sum_j\frac{\partial\varphi_i}{\partial
x_j}\dd x_j$，且 $\varphi^*(a\,\dd y_{i_1}\wedge\dots\wedge
\dd y_{i_k}) =
(a\circ\varphi)\,\dd\varphi_{i_1}\wedge\dots\wedge
\dd\varphi_{i_k}$。
\end{definition}

\begin{theorem}\label{thm:b3:forms:pullback}
(a) $\varphi^*(\omega \wedge \eta) = \varphi^*\omega \wedge
\varphi^*\eta$ 和 $(\psi \circ \varphi)^* = \varphi^* \circ
\psi^*$。
(b) $\varphi^*(\dd\omega) = \dd(\varphi^*\omega)$：外微分
与每一个光滑替换可交换，
--- 正是这一恒等式使它成为本理论中\emph{真正的}
导数。
(c) 若 $\varphi\colon U \to V$ 是
$\R^n$ 中开集之间的光滑映射，且 $\omega = a\,\dd y_1 \wedge \dots \wedge \dd
y_n$，则 $\varphi^*\omega = (a \circ
\varphi)\,\det\bigl(D\varphi\bigr)\,\dd x_1 \wedge \dots
\wedge \dd x_n$。
\end{theorem}

\begin{proof}
(a) 关于线性拉回的逐点陈述
（见\cref{prop:b3:forms:linearpullback}），再加上链式法则
$D(\psi\circ\varphi)(x) = D\psi(\varphi(x))\,D\varphi(x)$。
(b) 对 $0$-形式 $f$：$\varphi^*(\dd f) = \dd f \circ
D\varphi = \dd(f \circ \varphi)$ 即链式法则。对
$\omega = a\,\dd y_I$，利用 (a)，$\varphi^*\omega =
(a\circ\varphi)\,\dd\varphi_{i_1}\wedge\dots\wedge
\dd\varphi_{i_k}$，于是由莱布尼茨法则
（见\cref{thm:b3:forms:dsquare}(a)）和 $\dd\dd\varphi_{i_r} =
0$：
\[
\dd(\varphi^*\omega) = \dd(a\circ\varphi) \wedge
\dd\varphi_{i_1}\wedge\dots\wedge\dd\varphi_{i_k}
= \varphi^*(\dd a) \wedge \varphi^*(\dd y_I)
= \varphi^*(\dd a \wedge \dd y_I) =
\varphi^*(\dd\omega).
\]
(c) 逐点地，这恰好就是最高次形式上的 $u^*\alpha = (\det u)\alpha$
（见\cref{prop:b3:forms:linearpullback}(1)
，取 $u = D\varphi(x)$）。
\end{proof}

\begin{example}[极坐标]\label{ex:b3:forms:polar}
设 $\varphi(r, \theta) = (r\cos\theta, r\sin\theta)$：
$\varphi^*\dd x = \cos\theta\,\dd r -
r\sin\theta\,\dd\theta$，$\varphi^*\dd y = \sin\theta\,\dd r
+ r\cos\theta\,\dd\theta$，于是
\[
\varphi^*(\dd x \wedge \dd y) =
(\cos\theta\,\dd r - r\sin\theta\,\dd\theta) \wedge
(\sin\theta\,\dd r + r\cos\theta\,\dd\theta)
= r\,\dd r \wedge \dd\theta,
\]
\cref{ex:b3:product:polar}中的雅可比行列式借纯代数
自然出现 --- 无需任何测度论。\cref{exo:b3:forms:9}把这一注记化为
一个命题：对于定向积分，换元公式
\emph{正是}拉回公式。
\end{example}

\section{闭形式与恰当形式；庞加莱引理}

\begin{definition}\label{def:b3:forms:closedexact}
$\omega \in \Omega^k(U)$ 称为\emph{闭形式}\index{闭形式}，若
$\dd\omega = 0$；称为\emph{恰当形式}\index{恰当形式}，若
$\omega = \dd\eta$ 对某个 $\eta \in \Omega^{k-1}(U)$ 成立（称为
\emph{原形式}，其外微分即为 $\omega$）。恰当形式 $\Rightarrow$ 闭形式，由
$\dd^2 = 0$；其逆命题是关于
\emph{形状}的问题，即 $U$ 具有怎样的形状。
\end{definition}

\begin{example}[角形式]\label{ex:b3:forms:angular}
在 $U = \R^2 \setminus \{0\}$ 上，
\[
\omega_\theta = \frac{x\,\dd y - y\,\dd x}{x^2 + y^2}
\]
是闭形式（直接计算，见
\cref{exo:b3:forms:4}），但不是恰当形式：它沿单位圆的
积分为 $2\pi \neq 0$，而恰当形式沿闭曲线的
积分恒为零（见\cref{prop:b3:forms:exactloop}）。
局部地，$\omega_\theta = \dd\theta$ 对极角的任一光滑选取
$\theta$ 成立 --- 由此得名，
并揭示障碍：在整个
$U$ 上不存在这样的选取。这一形式串起环绕数
（见\cref{sec:b3:forms:winding}），并借此通向
\cref{ch:b3:residues}中的留数定理。
\end{example}

\begin{theorem}[庞加莱引理]\label{thm:b3:forms:poincare}
设 $U \subseteq \R^n$ 为开集，且为\emph{星形}的（关于
$0$）。每个闭 $k$-形式（定义在 $U$ 上，$k \geq 1$）皆
为恰当形式。\index{庞加莱引理}
\end{theorem}

\begin{proof}
构造一个线性\emph{同伦算子} $h\colon
\Omega^k(U) \to \Omega^{k-1}(U)$，使得
\begin{equation}\label{eq:b3:forms:homotopy}
\dd(h\omega) + h(\dd\omega) = \omega
\qquad (k \geq 1);
\end{equation}
若 $\dd\omega = 0$，则 $\omega = \dd(h\omega)$，于是
即得证。对 $\omega = \sum_I a_I\,\dd x_I$，令
\[
h\omega = \sum_{\abs I = k}\ \sum_{r=1}^{k}(-1)^{r-1}
\Bigl(\int_0^1 t^{k-1}a_I(tx)\,\dd t\Bigr)\,
x_{i_r}\,\dd x_{i_1}\wedge\dots\wedge
\widehat{\dd x_{i_r}}\wedge\dots\wedge\dd x_{i_k}
\]
（帽号表示删去一个因子；积分关于 $x$ 光滑，
由积分号下求导得出，
见\cref{thm:b3:lebesgue:paramdiff}，所有导数
在紧集上被控制）。验证
\eqref{eq:b3:forms:homotopy} 是一项一生只做一次的计算，
故完整写出。固定 $I$，取 $\omega =
a\,\dd x_I$（由线性性）。首先，
\[
\dd(h\omega) = k\Bigl(\int_0^1t^{k-1}a(tx)\dd t\Bigr)\dd x_I
+ \sum_{r=1}^k(-1)^{r-1}\sum_{j=1}^n
\Bigl(\int_0^1t^{k}\,\partial_ja(tx)\,\dd t\Bigr)
x_{i_r}\,\dd x_j\wedge\dd x_{I\setminus i_r} :
\]
第一组收集的是 $\dd$ 作用于
因子 $x_{i_r}$ 的那些项 --- 楔积 $\dd x_{i_r}\wedge\dd
x_{I\setminus i_r}$ 重新拼成 $\dd x_I$，带有符号
$(-1)^{r-1}$ 恰好抵消前置系数，而 $k$ 个
$r$ 的取值给出因子 $k$ --- 而第二组
收集 $\dd$ 作用于积分号的那些项（链式法则
提出 $t\,\partial_ja(tx)$）。接着，$\dd\omega =
\sum_j\partial_ja\,\dd x_j\wedge\dd x_I$，并应用
$h$ 在 $k+1$ 次情形的定义，指标 $j$ 占据
第一个位置：
\[
h(\dd\omega) = \sum_{j=1}^n\Bigl(\int_0^1t^{k}
\partial_ja(tx)\dd t\Bigr)x_j\,\dd x_I
- \sum_{j=1}^n\sum_{r=1}^k(-1)^{r-1}
\Bigl(\int_0^1t^{k}\partial_ja(tx)\dd t\Bigr)
x_{i_r}\,\dd x_j\wedge\dd x_{I\setminus i_r} .
\]
两个二重和在 $\dd(h\omega) + h(\dd\omega)$ 中相消，
故其等于
\[
\Bigl(\int_0^1\bigl(k\,t^{k-1}a(tx) +
t^k{\textstyle\sum_j}x_j\,\partial_ja(tx)\bigr)\dd t\Bigr)
\dd x_I
= \Bigl(\int_0^1\frac{\dd}{\dd t}\bigl(t^ka(tx)\bigr)\dd
t\Bigr)\dd x_I
= a(x)\,\dd x_I = \omega,
\]
由微积分基本定理。星形性
在必要之处起了作用：$tx \in U$ 对 $t \in \intcc01$ 成立，
从而 $a(tx)$ 有意义。
\end{proof}

\begin{remark}
当 $k = 1$ 且 $\omega = \sum_j a_j\dd x_j$ 时，其原函数
为 $f(x) = \int_0^1\sum_ja_j(tx)\,x_j\,\dd t$ --- 即
$\omega$ 沿线段 $[0, x]$ 的线积分。本定理即
本科二年级多元情形的结论“雅可比矩阵对称的场
是梯度场”，如今推广到所有次数。角
形式（\cref{ex:b3:forms:angular}）表明关于
$U$ 的假设并非可有可无：$\R^2\setminus\{0\}$ 不是
星形的，且在其上闭性并不蕴含恰当性。
在一般开集上保留的部分，由
\emph{德拉姆上同调} $H^k(U) =
\ker\dd/\operatorname{im}\dd$ 度量 ---
见\cref{exo:b3:forms:12}，给出第一个非平凡的计算。
\end{remark}

\section{子流形上的定向与积分}

对 $k$-形式作积分，需要 $k$ 维有向
区域。回忆在\cref{ch:b3:submanifolds}中
（\cref{thm:b3:submanifolds:characterizations}）可知，一个
$k$ 维子流形 $M \subseteq \R^n$ 局部地是
正则参数化 $\gamma\colon V \to M \cap W$（$V
\subseteq \R^k$ 为开集，$\gamma$ 为到其像的同胚，
且其微分为单射）。

\begin{definition}\label{def:b3:forms:orientation}
称 $M$ 的\emph{定向}\index{定向}为这样一种选择：
对每个 $p \in M$，在切空间 $T_pM$ 的两类定向基中选定一类，
且该选择是\emph{局部
相容的}：每一点附近都存在一个参数化
$\gamma$，其坐标标架 $(\partial_1\gamma, \dots,
\partial_k\gamma)$ 在其定义域的每一点处都是正定向的。
这样的参数化称为\emph{正向的}。
若定向存在，则称 $M$ 为\emph{可定向的}；
莫比乌斯带表明此性质可能失效。本章中的子流形
均已定向。
\end{definition}

\begin{definition}[形式的积分]\label{def:b3:forms:integral}
设 $M$ 为有定向的 $k$ 维子流形，$\omega$ 为
$k$-形式，它在 $M$ 的某邻域上有定义，并且
$\operatorname{supp}\omega \cap M$ 为紧集。
(a) 若 $\operatorname{supp}\omega \cap M \subseteq
\gamma(V)$（对某个单一的正向参数化），则令
\[
\int_M\omega = \int_V\gamma^*\omega
\]
--- 右边是 $V$ 上的勒贝格积分
（见\cref{ch:b3:product}）中被积函数的系数 $g$，
即 $\gamma^*\omega = g\,\dd u_1\wedge\dots\wedge\dd u_k$，它
是连续的，且具有紧支集。
(b) 一般地，选取有限个正向参数化
$\gamma_i(V_i)$ 覆盖紧集
$\operatorname{supp}\omega \cap M$，并取从属的
单位分解 $(\chi_i)$
（\cref{lem:b3:forms:partition}），并令 $\int_M\omega =
\sum_i\int_M\chi_i\,\omega$，每一项由 (a) 计算。
对 $k = 1$ 的曲线，若它由
$\gamma\colon\intcc ab\to\R^n$ 参数化，则记
$\int_\gamma\omega = \int_a^b\gamma^*\omega$，
且不要求单射。
\end{definition}

\begin{lemma}[相容性]\label{lem:b3:forms:consistency}
定义 (a) 与所取的正向参数化无关，
定义 (b) 既与覆盖无关，
也与单位分解无关。此外，若
$\Phi$ 是两个有定向子流形邻域之间的微分同胚，
使 $\Phi(M) = M'$，并把
$M$ 的正向标架映为 $M'$ 的正向标架，且在
$M$ 的每个连通分支上的某点处，则 $\int_{M'}\omega =
\int_M\Phi^*\omega$.
\end{lemma}

\begin{proof}
(a) 设 $\gamma\colon V \to M$，$\delta\colon V' \to M$ 为
正向参数化，其像包含
$\operatorname{supp}\omega\cap M$。过渡映射 $\tau =
\delta^{-1}\circ\gamma$ 是 $V, V'$ 中相关开子集之间的
微分同胚（光滑性见
\cref{thm:b3:submanifolds:characterizations}，借助局部图像
描述），且在那里 $\gamma = \delta\circ\tau$，因此
$\gamma^*\omega = \tau^*(\delta^*\omega)$
（\cref{thm:b3:forms:pullback}(a)）。记 $\delta^*\omega =
g\,\dd u_1\wedge\dots\wedge\dd u_k$；于是
（\cref{thm:b3:forms:pullback}(c)）得 $\tau^*(\delta^*\omega) =
(g\circ\tau)\,\det(D\tau)\,\dd u_1\wedge\dots\wedge\dd u_k$.
两个标架都是正向的，故 $D\tau$ 把正基映为
正基：$\det D\tau > 0$，因此 $\det D\tau =
\abs{\det D\tau}$，并由变量替换定理
（\cref{thm:b3:product:changeofvar}）得
$\int(g\circ\tau)\abs{\det D\tau} = \int g$：两个
积分相等。\emph{这正是定向存在的
根本原因}：若无符号控制，雅可比行列式及其
绝对值不同，积分便无定义。
(b) 若 $(\chi_i)$ 与 $(\tilde\chi_j)$ 为两个可容许的
单位分解（包括覆盖在内），则由 (a) 及有限
可加性，$\sum_i\int\chi_i\omega =
\sum_{i,j}\int\chi_i\tilde\chi_j\omega =
\sum_j\int\tilde\chi_j\omega$，每个双重项
可在任一坐标卡中计算。最后一句：若 $\gamma$ 取遍
$M$ 的正向参数化，则 $\Phi\circ\gamma$ 取遍
$M'$ 的正向参数化（定向的
比较是局部常值的，并在每个连通分支上于
一点处固定），且 $(\Phi\circ\gamma)^*\omega =
\gamma^*(\Phi^*\omega)$.
\end{proof}

\begin{lemma}[单位分解，紧情形]
\label{lem:b3:forms:partition}
设 $K \subseteq \R^n$ 为紧集，$W_1, \dots, W_m$ 为
覆盖 $K$ 的开集。则存在 $\chi_1, \dots, \chi_m \in
\mathcal C^\infty(\R^n)$，满足 $0 \leq \chi_i \leq 1$，
$\operatorname{supp}\chi_i \subseteq W_i$ 为紧集，并且
$\sum\chi_i = 1$ 在
$K$ 的某一邻域上成立。\index{单位分解}
\end{lemma}

\begin{proof}
每个 $x \in K$ 都落在某个 $W_{i(x)}$ 中，且存在闭球
$\bar B(x, 2r_x) \subseteq W_{i(x)}$，由紧性可取出
$x_1, \dots, x_N$，使球 $B(x_s, r_{x_s})$ 覆盖
$K$。对每个 $s$，取光滑隆起函数 $\theta_s \in \mathcal
C^\infty$，$0 \leq \theta_s \leq 1$，$\theta_s = 1$ 在
$\bar B(x_s, r_{x_s})$ 上，$\operatorname{supp}\theta_s
\subseteq B(x_s, 2r_{x_s})$（对
半径为 $\frac32r_{x_s}$ 的球之示性函数作磨光，
见\cref{thm:b3:lp:regularization}）。为每个 $s$ 指定一个
指标 $i(s)$，使 $B(x_s, 2r_{x_s}) \subseteq W_{i(s)}$，
并令 $\Theta_i = \sum_{i(s) = i}\theta_s$。
在开集 $\Omega_0 = \{\sum_j\Theta_j > \tfrac12\} \supseteq K$ 上，
函数 $\Theta_i/\sum_j\Theta_j$ 即满足要求，但仅在该集上有定义；
为使之整体化，取 $\rho \in \mathcal
C^\infty(\R^n)$，使 $\rho = 0$，在 $\sum_j\Theta_j
\geq 1$ 处，$\rho > 0$，在 $\sum_j\Theta_j \leq \tfrac12$ 处
（对 $1 - \sum_j\Theta_j$ 的适当截断函数作磨光），
并令
\[
\chi_i = \frac{\Theta_i}{\rho + \sum_j\Theta_j} .
\]
分母处处 $> 0$，且等于
$\sum_j\Theta_j$，而 $\{\sum_j\Theta_j \geq 1\}$ 是
$K$ 的开邻域（$K$ 的每一点都落在某个球
$B(x_s, r_{x_s})$ 中，而在其上 $\theta_s = 1$）；在此邻域上
$\sum_i\chi_i = 1$。支集与界显然成立。
\end{proof}

\section{斯托克斯定理}

\begin{definition}\label{def:b3:forms:boundary}
\emph{带边 $k$ 维子流形}\index{带边子流形}
$M \subseteq \R^n$ 是由正则参数化覆盖的集合，这些参数化分两类：
\emph{内部坐标卡}，
$\gamma\colon V \to M \cap W$（$V \subseteq \R^k$ 为开集），
以及 \emph{边界坐标卡} $\gamma\colon V \cap H^k \to M
\cap W$，其中 $H^k = \{u \in \R^k : u_k \geq 0\}$，
且 $\gamma$ 可光滑且正则地延拓到开集 $V$。
\emph{边界} $\partial M$ 是指
在 $u_k = 0$ 处取到的点之集合；它是一个
不带边的 $(k-1)$ 维子流形，由映射 $u' \mapsto \gamma(u',
0)$ 参数化。$M$ 的一个定向
通过\emph{外法线优先}规则诱导 $\partial M$ 的定向：对 $p \in \partial M$，
$(w_1, \dots, w_{k-1})$ 是 $T_p\partial M$ 的一组基，
它为正，当且仅当 $(\nu, w_1, \dots, w_{k-1})$ 是
$T_pM$ 的正基，其中 $\nu \in T_pM \setminus T_p\partial M$，
指向 $M$ 外部（在边界坐标卡中：$\nu =
-\partial_k\gamma$，允许相差切向分量 ---
定向类不依赖这些分量）。
\end{definition}

\begin{omfigure}
\begin{tikzpicture}[scale=1.05]
  \fill[omDef!10] (-2,0) arc (180:0:2) -- cycle;
  \draw[thick, omDef] (-2,0) arc (180:0:2);
  \draw[thick, omThm] (-2,0) -- (2,0);
  \node[omDef] at (0.1,1.35) {$M$};
  \node[omThm, below] at (1.55,-0.05) {$\partial M$};
  \draw[->, thick, omExo] (1.415,1.415) -- (2.05,2.05);
  \node[omExo, right] at (2.02,2.12) {$\nu$};
  \draw[->, thick] (1.29,1.53) arc (50:75:2);
  \draw[->, thick, omExo] (-0.6,0) -- (-0.6,-0.65);
  \node[omExo, below] at (-0.6,-0.65) {$\nu$};
  \draw[->, thick] (-0.35,0) -- (0.55,0);
  \node[below] at (0.1,-0.14) {诱导};
\end{tikzpicture}

{\small 外法线优先规则：在每个边界点处，
将外向向量 $\nu$ 置于首位；
能与之补成 $M$ 的正标架的基，给出 $\partial M$ 的定向。
对具标准定向的平面区域，这正是
格林–黎曼的逆时针规则。}
\end{omfigure}

\begin{lemma}[半空间上的斯托克斯定理]
\label{lem:b3:forms:halfspace}
设 $\eta$ 为光滑的 $(k-1)$ 形式，它在 $\R^k$ 上
具紧支集，则
\[
\int_{H^k}\dd\eta = \int_{\partial H^k}\eta,
\]
其中 $H^k$ 带有 $\R^k$ 的标准定向，
$\partial H^k = \{u_k = 0\} \cong \R^{k-1}$ 上的诱导定向，
它等于 $(-1)^k$ 倍，
即 $\R^{k-1}$ 的标准定向。
\end{lemma}

\begin{proof}
先做定向核算：
沿 $\partial H^k$ 的外法向为 $-e_k$，并
\[
\det(-e_k, e_1, \dots, e_{k-1}) =
-\det(e_k, e_1, \dots, e_{k-1}) = -(-1)^{k-1}\det(e_1,
\dots, e_k) = (-1)^k :
\]
标架 $(e_1, \dots, e_{k-1})$ 是 $\partial H^k$ 的一组基，
它关于诱导定向为正，当且仅当 $k$
为偶数，由此得所述比较。由线性性，取 $\eta
= f\,\dd u_1\wedge\dots\wedge\widehat{\dd
u_i}\wedge\dots\wedge\dd u_k$, $f \in \mathcal
C^\infty_c(\R^k)$，则 $\dd\eta =
(-1)^{i-1}\,\partial_if\,\dd u_1\wedge\dots\wedge\dd u_k$
（把 $\dd u_i$ 移入其位置需 $i - 1$ 次对换）。两种
情形均由托内利--富比尼
（\cref{thm:b3:product:fubini}）及一元
微积分基本定理推出。

\emph{情形 $i < k$。}先对 $u_i$ 在 $\R$ 上积分：
$\int_\R\partial_if\,\dd u_i = 0$（由紧支集），
故 $\int_{H^k}\dd\eta = 0$。又 $\eta$ 在
$\{u_k = 0\}$ 上的限制含因子 $\dd u_k$，
它限制为 $0$（此处 $u_k$ 为常数）：$\int_{\partial H^k}\eta =
0$ 亦成立。

\emph{情形 $i = k$。}先对 $u_k$
在 $\intco0\infty$ 上积分：
\begin{align*}
\int_{H^k}\dd\eta &= (-1)^{k-1}\int_{\R^{k-1}}
\Bigl(\int_0^\infty\partial_kf\,\dd u_k\Bigr)\dd u' \\
&= (-1)^{k-1}\int_{\R^{k-1}}\bigl(0 - f(u', 0)\bigr)\dd u'
= (-1)^k\int_{\R^{k-1}}f(u',0)\,\dd u' .
\end{align*}
边界一侧，$\eta$ 的限制为 $f(u',
0)\,\dd u_1\wedge\dots\wedge\dd u_{k-1}$，诱导
定向为标准定向的 $(-1)^k$ 倍，于是
$\int_{\partial H^k}\eta =
(-1)^k\int_{\R^{k-1}}f(u',0)\,\dd u'$。两侧一致。
\end{proof}

\begin{theorem}[斯托克斯]\label{thm:b3:forms:stokes}
设 $M \subseteq \R^n$ 为紧的定向
带边 $k$ 维子流形，$\partial M$ 赋予诱导定向，
并设 $\omega$ 为光滑的 $(k-1)$ 形式，定义在
$M$ 的邻域上。则\index{斯托克斯定理}
\[
\int_M\dd\omega = \int_{\partial M}\omega .
\]
特别地，若 $\partial M = \varnothing$，
则 $\int_M\dd\omega = 0$。
\end{theorem}

\begin{proof}
用有限个直接坐标卡（内部或边界）的像覆盖紧集 $M$，
并取单位分解
$(\chi_i)$，从属于对应的开集 $W_i$
（$\R^n$ 中，见\cref{lem:b3:forms:partition}），且 $\sum\chi_i
= 1$ 于 $M$ 的某邻域上。在该邻域上
$\dd(\sum\chi_i) = 0$，于是
\[
\int_M\dd\omega = \sum_i\int_M\dd(\chi_i\omega),
\qquad
\int_{\partial M}\omega =
\sum_i\int_{\partial M}\chi_i\omega :
\]
两边都可加，故只需对
支集位于单个坐标卡像中的形式证明该定理。

\emph{内部坐标卡。}若 $\operatorname{supp}\omega \cap M
\subseteq \gamma(V)$，$V$ 为 $\R^k$ 中的开集：令 $\eta =
\gamma^*\omega$ 以零延拓到 $\R^k$（光滑，紧支集
位于 $V$ 中），并应用\cref{lem:b3:forms:halfspace}，
使支集远离 $\partial H^k$（将 $V$ 平移入
开上半空间 --- 或直接重复 $i < k$ 情形的
在整个 $\R^k$ 上的计算）：$\int_M\dd\omega =
\int_{\R^k}\dd\eta = 0$，又 $\int_{\partial M}\omega = 0$，
因为 $\omega$ 在 $\partial M$ 附近为零。

\emph{边界坐标卡。}若 $\operatorname{supp}\omega\cap M
\subseteq \gamma(V \cap H^k)$：令 $\eta = \gamma^*\omega$
以零延拓，
$\gamma^*(\dd\omega) = \dd\eta$
（\cref{thm:b3:forms:pullback}(b)），于是
由\cref{def:b3:forms:integral}，有：
\[
\int_M\dd\omega = \int_{H^k}\dd\eta
\overset{\text{\cref{lem:b3:forms:halfspace}}}{=}
\int_{\partial H^k}\eta .
\]
还需将右端与 $\int_{\partial
M}\omega$ 相等。边界参数化为 $\beta(u') =
\gamma(u', 0)$，而 $\beta^*\omega$ 为
$\eta$ 在 $\{u_k = 0\}$ 上的限制（沿包含映射 $u'
\mapsto (u', 0)$ 与 $\gamma$ 复合的拉回）。定向
比较在两侧都是 $(-1)^k$：标架
$(\partial_1\beta, \dots, \partial_{k-1}\beta)$ 位于
$\partial M$ 的诱导定向中，符号为
$\det(-e_k, e_1, \dots, e_{k-1}) = (-1)^k$，在坐标卡中算得
（外法向量拉回为 $-e_k$），恰是
联系 $\partial H^k$ 的诱导定向
与 $\R^{k-1}$ 标准定向的符号
（见\cref{lem:b3:forms:halfspace}）。两个符号约定
相互抵消：
$\int_{\partial H^k}\eta = \int_{\partial M}\omega$。
\end{proof}

\begin{example}[经典定理]
\label{ex:b3:forms:classical}
设 $D \subseteq \R^2$ 为具光滑边界的紧区域，
取标准定向。对 $\omega = P\,\dd x +
Q\,\dd y$：$\dd\omega = \bigl(\partial_xQ -
\partial_yP\bigr)\dd x\wedge\dd y$，斯托克斯公式为
\[
\int_D\Bigl(\frac{\partial Q}{\partial x} -
\frac{\partial P}{\partial y}\Bigr)\dd x\,\dd y
= \oint_{\partial D}P\,\dd x + Q\,\dd y :
\]
\emph{格林--黎曼公式}，已在
本科二年级分册中对初等区域证明，并在此以自然的一般形式给出。在 $\R^3$ 中，
将斯托克斯公式用于向量场通量的 $2$-形式，
在紧区域上得到\emph{散度定理}：
$\int_\Omega\operatorname{div}F =
\int_{\partial\Omega}\langle F, \nu\rangle\,\dd S$，又将
它用于带边曲面上的 $1$-形式，则得到
经典的\emph{开尔文--斯托克斯}旋度定理；
详见\cref{exo:b3:forms:7}，其中列出两种对应关系。
\end{example}

\section{环绕数}\label{sec:b3:forms:winding}

\begin{definition}\label{def:b3:forms:winding}
设 $\gamma\colon\intcc01\to\R^2\setminus\{a\}$ 为光滑
闭曲线。其\emph{环绕数}\index{环绕
数}关于点 $a$ 定义为：
\[
\operatorname{Ind}_\gamma(a) =
\frac1{2\pi}\int_\gamma\omega_\theta^a,
\qquad
\omega_\theta^a = \frac{(x - a_1)\,\dd y - (y - a_2)\,\dd
x}{(x - a_1)^2 + (y - a_2)^2} .
\]
\end{definition}

\begin{proposition}\label{prop:b3:forms:windinginteger}
$\operatorname{Ind}_\gamma(a) \in \Z$；作为 $a$ 的函数，
它在 $\R^2\setminus\gamma(\intcc01)$ 的每个连通分支上为常数，
并在无界分支上为零。
对于 $\gamma(t) = a + r(\cos2\pi t, \sin2\pi t)$，有
$\operatorname{Ind}_\gamma(a) = 1$。
\end{proposition}

\begin{proof}
取 $a = 0$，记 $\gamma = (x, y)$，$\rho =
\norm\gamma > 0$，$\vartheta(t) =
\int_0^t\gamma^*\omega_\theta$，使得 $\vartheta' =
\frac{xy' - yx'}{x^2 + y^2}$。在复记号下令 $u(t)
= \gamma(t)\,\rho(t)^{-1}\eu^{-\iu\vartheta(t)}$；则
$\abs u = 1$，且直接计算可得
\[
\frac{u'}{u} = \frac{\gamma'}{\gamma} - \frac{\rho'}{\rho}
- \iu\vartheta'
= \frac{\bar\gamma\gamma'}{\abs\gamma^2} -
\frac{\rho'}{\rho} - \iu\,\frac{xy' -
yx'}{\rho^2}
= \frac{(xx' + yy')}{\rho^2} - \frac{\rho'}{\rho} = 0,
\]
由于 $\bar\gamma\gamma' = (xx' + yy') + \iu(xy' - yx')$，且
$\rho\rho' = xx' + yy'$。故 $u$ 为常数：$\gamma(t) =
\rho(t)\,c\,\eu^{\iu\vartheta(t)}$，其中 $\abs c = 1$，且
$\gamma(1) = \gamma(0)$ 及 $\rho(1) = \rho(0)$ 蕴含
$\eu^{\iu\vartheta(1)} = \eu^{\iu\vartheta(0)} = 1$:
$\vartheta(1) \in 2\pi\Z$，即
$\operatorname{Ind}_\gamma(0) \in \Z$。作为 $a$ 的函数，
在紧曲线的开补集上，定义式中的积分
是连续的（由\cref{thm:b3:lebesgue:paramcont}，
在每个 $a$ 的邻域上有可积控制）；
连续的整值函数局部为常数，因而
在各连通分支上为常数。当 $\norm a$ 很大时，被积函数
$O(1/\norm a)$ 关于 $t$ 一致成立，于是环绕数趋于 $0$，
从而在无界分支上为零。对于圆周，
直接有 $\gamma^*\omega_\theta^a = 2\pi\,\dd t$。
\end{proof}

\begin{remark}
在 $\C \cong \R^2$ 下，$\frac{\dd z}{z} = \frac{x\,\dd x +
y\,\dd y}{x^2 + y^2} + \iu\,\omega_\theta$，于是
$\operatorname{Ind}_\gamma(a) =
\frac1{2\iu\pi}\oint_\gamma\frac{\dd z}{z - a}$：这正是
\cref{ch:b3:residues}中的环绕数，并且
\cref{prop:b3:forms:windinginteger}以实变量方法重新证明了它的
整值性与局部常数性 ---
这是留数定理中的拓扑一半，如今建立在
斯托克斯公式之上。
\end{remark}

\begin{proposition}\label{prop:b3:forms:exactloop}
若 $\omega = \dd f$ 在开集 $U$ 上恰当，且
$\gamma\colon\intcc01\to U$ 为闭曲线，则
$\int_\gamma\omega = 0$。
\end{proposition}

\begin{proof}
$\int_\gamma\dd f = \int_0^1(f\circ\gamma)'(t)\,\dd t =
f(\gamma(1)) - f(\gamma(0)) = 0$ --- 链式法则
将 $\gamma^*(\dd f)$ 等同于 $(f\circ\gamma)'\,\dd t$。
\end{proof}

\begin{method}\label{met:b3:forms:method}
计算微分形式时：（1）机械化 --- 楔积按符号重排，
$\dd$ 对系数求导并引入新的 $\dd
x_j$，拉回时作代换；信任代数运算，它
编码了每一个雅可比行列式。（2）要在子流形上对形式积分：
直接参数化，作拉回，再对系数积分；
定向是唯一的陷阱 --- 检验一个标架。（3）要证明积分恒等式，
就寻找斯托克斯型结构：被积函数是否恰当？
区域是否为边界？（4）要比较两个“平行”子流形上的积分，
只需对二者之间的区域应用斯托克斯公式
（形变论证，
见\cref{exo:b3:forms:10}）。（5）闭形式的非零积分
证明存在拓扑障碍 --- 无原函数，
无收缩，无零点延拓：这正是
周末问题排除球体收缩映射的方法。
\end{method}

\section{练习}

\begin{exercise}[$\star$]\label{exo:b3:forms:1}
在 $\R^3$ 上，设 $\omega = x\,\dd y - z\,\dd x$ 及 $\eta =
\dd x + y\,\dd z$。计算 $\omega\wedge\eta$，$\dd\omega$，
$\dd\eta$，$\dd(\omega\wedge\eta)$，并在此例上验证
分次莱布尼茨法则。
\end{exercise}

\begin{exercise}[$\star$]\label{exo:b3:forms:2}
在 $\R^3$ 上辨认 $\dd$ 的三种形式：
对 $f \in \Omega^0$，$\dd f \leftrightarrow \nabla f$；对
功形式 $\omega_F = F_1\dd x + F_2\dd y + F_3\dd z$，
$\dd\omega_F \leftrightarrow \operatorname{curl}F$；对
通量形式 $\sigma_F = F_1\,\dd y\wedge\dd z + F_2\,\dd
z\wedge\dd x + F_3\,\dd x\wedge\dd y$，$\dd\sigma_F
\leftrightarrow \operatorname{div}F$。由 $\dd^2 =
0$ 推出恒等式 $\operatorname{curl}\nabla f = 0$ 与
$\operatorname{div}\operatorname{curl}F = 0$。
\end{exercise}

\begin{exercise}[$\star\star$]\label{exo:b3:forms:3}
判断下列每个 $1$-形式是否为闭形式，在其定义域上是否恰当，
并在存在时求出原函数：
(a) $(2xy + z^2)\,\dd x + x^2\,\dd y + 2xz\,\dd z$ 在
$\R^3$ 上；
(b) $\dfrac{x\,\dd x + y\,\dd y}{x^2 + y^2}$ 在
$\R^2\setminus\{0\}$ 上；
(c) $\dfrac{-y\,\dd x + x\,\dd y}{x^2 + y^2}$ 在
半平面 $\{x > 0\}$ 上。
\end{exercise}

\begin{exercise}[$\star\star$]\label{exo:b3:forms:4}
（角形式）验证 $\omega_\theta$
（见\cref{ex:b3:forms:angular}）为闭形式；计算
$\int_\gamma\omega_\theta$，其中 $\gamma$ 是半径
$r$、以 $0$ 为中心的圆周；由此断言 $\omega_\theta$ 在
$\R^2\setminus\{0\}$ 上不恰当，并说明 $\R^2\setminus\{0\}$
关于其中任何一点都不是星形的（两种途径：
经由\cref{thm:b3:forms:poincare}，
或直接由几何得出）。
\end{exercise}

\begin{exercise}[$\star\star$]\label{exo:b3:forms:5}
（超曲面的面积形式）设 $M \subseteq \R^n$ 为
紧的可定向超曲面，其定向由
单位法向量场 $\nu$ 给出（$(w_1, \dots, w_{n-1})$ 为正向，
当且仅当 $(\nu, w_1, \dots, w_{n-1})$ 在 $\R^n$ 中为正向）。证明
$(n-1)$ 形式 $\sigma_\nu =
\sum_i(-1)^{i-1}\nu_i\,\dd x_1\wedge\dots\wedge\widehat{\dd
x_i}\wedge\dots\wedge\dd x_n$ 在 $M$ 上的限制为
\emph{面积形式}：对正向参数化 $\gamma$，
$\gamma^*\sigma_\nu = \sqrt{\det G}\,\dd u_1\wedge\dots
\wedge\dd u_{n-1}$，其中 $G = ({}^t D\gamma)(D\gamma)$ 为
格拉姆矩阵。\emph{（注意
$(\gamma^*\sigma_\nu)(e_1, \dots, e_{n-1}) = \det(\nu,
\partial_1\gamma, \dots, \partial_{n-1}\gamma)$，并将
此行列式平方。）} 计算 $\int_{S^2}\sigma_\nu = 4\pi$。
\end{exercise}

\begin{exercise}[$\star\star$]\label{exo:b3:forms:6}
计算 $\int_{S^2}\omega$，其中 $\omega = x\,\dd y\wedge\dd z
+ y\,\dd z\wedge\dd x + z\,\dd x\wedge\dd y$：(a) 直接
用球坐标计算；(b) 借助单位球上的斯托克斯公式。
推出 $\operatorname{vol}(B^3) = \frac{4\pi}3$，
依据为 $S^2$ 的面积，并推广为
$n\operatorname{vol}(B^n) = \operatorname{area}(S^{n-1})$,
与\cref{thm:b3:product:ballvolume}一致。
\end{exercise}

\begin{exercise}[$\star\star$]\label{exo:b3:forms:7}
（“词典”）仔细地由
\cref{thm:b3:forms:stokes} 出发推导：(a) 在
$\R^3$ 中的散度定理（结合\cref{exo:b3:forms:2}与
\cref{exo:b3:forms:5}）；(b) 开尔文--斯托克斯定理
$\int_S\langle\operatorname{curl}F, \nu\rangle\,\dd S =
\oint_{\partial S}\langle F, \tau\rangle\,\dd\ell$ 对
$\R^3$ 中紧致的带边可定向曲面成立。检验
两种定向约定在上半球面上是否一致，
该上半球面以赤道为边界。
\end{exercise}

\begin{exercise}[$\star\star$]\label{exo:b3:forms:8}
（格林恒等式）设 $u, v$ 在
紧区域 $\Omega \subseteq \R^n$ 的邻域上光滑，且该区域
边界光滑。证明：
\[
\int_\Omega\bigl(u\,\Delta v + \langle\nabla u, \nabla
v\rangle\bigr) = \int_{\partial\Omega}u\,\partial_\nu
v\,\dd S,
\qquad
\int_\Omega\bigl(u\,\Delta v - v\,\Delta u\bigr) =
\int_{\partial\Omega}\bigl(u\,\partial_\nu v -
v\,\partial_\nu u\bigr)\dd S .
\]
由此推得：$\Omega$ 上的调和函数若在
$\partial\Omega$ 上为零，则恒为零，且两个调和函数
若边界值相同，则二者恒等 ---
即狄利克雷问题的唯一性，
见\cref{ch:b3:conformal}。
\end{exercise}

\begin{exercise}[$\star\star$]\label{exo:b3:forms:9}
（有向形式下的变量替换）设 $\varphi\colon U
\to V$ 为 $\R^n$ 中开集之间的微分同胚，且 $\det
D\varphi > 0$，$f$ 为连续函数，
其在 $V$ 中具有紧支集。证明拉回恒等式
$\int_U\varphi^*(f\,\dd x_1\wedge\dots\wedge\dd x_n) =
\int_Vf\,\dd x_1\wedge\dots\wedge\dd x_n$ 等价于
变量替换定理
（见\cref{thm:b3:product:changeofvar}）对此类 $\varphi$ 成立，并
精确说明雅可比行列式上的绝对值
去了哪里。
\end{exercise}

\begin{exercise}[$\star\star\star$]\label{exo:b3:forms:10}
（形变）设 $\omega$ 为\emph{闭} $2$-形式，
定义于 $\R^3\setminus\{0\}$ 上，$S_r$ 为半径为 $r$
以 $0$ 为中心的球面。证明 $\int_{S_r}\omega$
与 $r > 0$ 无关\emph{（对两半径之间的壳层应用斯托克斯公式，
注意两个边界的定向）}。应用于
\emph{立体角形式}
\[
\omega = \frac{x\,\dd y\wedge\dd z + y\,\dd z\wedge\dd x +
z\,\dd x\wedge\dd y}{(x^2 + y^2 + z^2)^{3/2}} :
\]
验证它是闭形式，计算 $\int_{S_r}\omega = 4\pi$，并
得出它是闭的，但在
$\R^3\setminus\{0\}$ 上不恰当 --- 它是
$\omega_\theta$ 的二维对应物，并且是高斯定律
在静电学中的几何内涵。
\end{exercise}

\begin{exercise}[$\star\star$]\label{exo:b3:forms:11}
设 $\gamma$ 为光滑闭曲线，位于
$\R^2\setminus\{0\}$ 中，且 $n = \operatorname{Ind}_\gamma(0)$。
证明 $\int_\gamma\omega = n\int_{S^1}\omega$，对
\emph{每个}闭 $1$-形式 $\omega$，
$\R^2\setminus\{0\}$ 上成立 \emph{（写作 $\omega = c\,\omega_\theta
+ \dd f$，由\cref{exo:b3:forms:12}得出）}。解释：在
去心平面上，环绕数是
使各周期为零的唯一障碍。
\end{exercise}

\begin{exercise}[$\star\star\star$]\label{exo:b3:forms:12}
（第一个德拉姆计算）证明每个闭 $1$-形式
$\omega$ 在 $U = \R^2\setminus\{0\}$ 上可唯一地写作
\[
\omega = c\,\omega_\theta + \dd f,
\qquad c = \frac1{2\pi}\int_{S^1}\omega,\quad f \in
\mathcal C^\infty(U) :
\]
定义 $f(p)$ 为 $\omega - c\,\omega_\theta$ 沿
从 $(1,0)$ 到 $p$ 的路径（先径向段，后圆
弧段）的积分，并证明结果与路径的选取无关，
恰因 $S^1$ 上的周期为零，并验证 $\dd
f = \omega - c\,\omega_\theta$。结论：
$H^1(\R^2\setminus\{0\}) \cong \R$，由
角形式生成。
\end{exercise}

\section{问题：布劳威尔不动点定理}

\begin{problem}[{周末问题 --- 无收缩，无
出路}]\label{pb:b3:forms:1}
布劳威尔定理断言：\emph{闭单位球
$\bar B = \bar B^n \subseteq \R^n$ 到自身的每个连续映射
都有一个不动点} --- 数学中的伟大定理之一，
其应用从博弈论（纳什
均衡）直至矩阵分析。微分形式证明是已知最简洁的：
斯托克斯公式表明球面不是球体的
收缩核，其余一切随之而来。全文中
$S = S^{n-1} = \partial\bar B$，$n \geq 2$，且
\[
\sigma = \sum_{i=1}^n(-1)^{i-1}x_i\,
\dd x_1\wedge\dots\wedge\widehat{\dd x_i}\wedge\dots\wedge
\dd x_n
\]
（帽号表示删去该因子）。

\medskip
\noindent\textbf{第 I 部分 --- 测量工具。}
\begin{enumerate}
  \item 计算 $\dd\sigma$，并由斯托克斯公式
        （见\cref{thm:b3:forms:stokes}）推出 $\int_S\sigma =
        n\operatorname{vol}(\bar B) > 0$，其中 $S$ 带有
        球的边界定向。
  \item 对 $n = 2$ 与 $n = 3$，将
        $\sigma$ 在 $S$ 上的限制与弧长形式和面积形式等同
        （见\cref{exo:b3:forms:5}，取 $\nu(x) = x$），然后
        直接重新计算 $\int_S\sigma$。
  \item 设 $W \subseteq \R^N$ 为开集，$\varphi\colon W
        \to \R^n$ 为光滑映射，且 $\norm{\varphi(x)} = 1$ 对
        所有 $x \in W$ 成立。证明
        $\varphi^*(\dd\sigma) = 0$。\emph{（对
        $\norm\varphi^2 = 1$ 求导：$D\varphi(x)$ 的像
        位于超平面 $\varphi(x)^\perp$ 中，维数为
        $n - 1$；然后应用
        \cref{prop:b3:forms:linearpullback}(2).)}
  \item 下面这段“证明”中错误在哪里？它断言
        $\int_S\sigma = 0$：“$S$ 紧且无
        边界，$\sigma$ 在 $S$ 上的限制是最高次
        形式，因而闭，故 $\int_S\sigma =
        \int_S\dd(\text{某物}) = 0$（由斯托克斯公式）”？
        （指出错误的那个词。）
  \item 用一段话说明第 II 部分的思路：
        被积的是什么、在何处积分，
        以及矛盾将从何处得出。
\end{enumerate}

\medskip
\noindent\textbf{第 II 部分 --- 不存在光滑收缩。} 用反证法，设
$r$ 为\emph{光滑收缩}，即
从球体到其球面的映射：$r$ 在
$\bar B$ 的某邻域上光滑，$r(\bar B) \subseteq S$，且 $r(x) = x$ 对所有
$x \in S$ 成立。
\begin{enumerate}[resume]
  \item 证明 $\int_Sr^*\sigma = \int_S\sigma$
        \emph{（在 $S$ 上，$r$ 限制为恒同：若
        $\gamma$ 为一段
        $S$ 的直接参数化，则 $r\circ\gamma = \gamma$）}。
  \item 在 $\bar B$ 上用斯托克斯公式，证明 $\int_Sr^*\sigma =
        \int_{\bar B}\dd(r^*\sigma)$.
  \item 证明 $\dd(r^*\sigma) = r^*(\dd\sigma) = 0$（
        将第 3 题应用于 $\varphi = r$），并得出结论：
        \emph{不存在 $\bar B \to S$ 的光滑收缩。}
  \item 手工处理被排除的情形 $n = 1$：
        直接证明不存在连续映射 $\intcc{-1}1 \to
        \{-1, 1\}$ 使两个端点都保持不动，并写出所用的
        定理的名称。
\end{enumerate}

\medskip
\noindent\textbf{第 III 部分 --- 光滑情形的布劳威尔定理。} 设 $g$
在 $\bar B$ 的某邻域上光滑，满足 $g(\bar B)
\subseteq \bar B$，且在 $\bar B$ 中\emph{没有}不动点。
\begin{enumerate}[resume]
  \item 证明 $\delta = \min_{x\in\bar B}\norm{g(x) - x} >
        0$.
  \item 对 $x \in \bar B$，令 $u(x) = \frac{x -
        g(x)}{\norm{x - g(x)}}$，并令 $r(x) = x +
        t(x)\,u(x)$ 为射线
        $\{x + tu(x) : t \geq 0\}$ 与 $S$ 的交点。解
        二次方程 $\norm{x + tu}^2 = 1$，得到
        \[
        t(x) = -\langle x, u(x)\rangle +
        \sqrt{1 - \norm x^2 + \langle x, u(x)\rangle^2}
        \;\geq\; 0 .
        \]
  \item 证明被开方式\emph{严格}为正：
        在 $\bar B$ 上，若 $1 - \norm x^2 + \langle x,
        u(x)\rangle^2 = 0$，则 $\norm x = 1$ 且 $\langle
        x, u(x)\rangle = 0$，即 $\langle x, x -
        g(x)\rangle = 0$，即 $\langle x, g(x)\rangle =
        1$；由柯西–施瓦茨不等式，取 $\norm x = 1$，
        $\norm{g(x)} \leq 1$，这迫使 $g(x) = x$ ---
        被排除。由此推出 $r$ 在
        $\bar B$ 的某邻域上光滑。
  \item 证明 $r(\bar B) \subseteq S$，且 $r(x) = x$ 对 $x
        \in S$ 成立\emph{（对 $\norm x = 1$，验证 $t(x) = 0$，
        利用 $\langle x, u(x)\rangle \geq 0$，后者又
        由 $\langle x, x - g(x)\rangle = 1 -
        \langle x, g(x)\rangle \geq 0$）}。结合第 II 部分，
        得出：\emph{$\bar B$ 的每个光滑自映射都有
        一个不动点。}
\end{enumerate}

\medskip
\noindent\textbf{第四部分 --- 连续布劳威尔定理。}设
$f\colon\bar B\to\bar B$ 为无不动点的连续映射。
\begin{enumerate}[resume]
  \item 证明 $\varepsilon = \min_{\bar B}\norm{f - x} > 0$，
        并构造一个\emph{多项式}映射 $p\colon\R^n\to
        \R^n$，使得 $\sup_{\bar B}\norm{p - f} <
        \varepsilon/2$（斯通–魏尔斯特拉斯，
        \cref{thm:b3:complete:stoneweierstrass}，逐分量
        进行——说明由标量逼近过渡到
        向量逼近的合理性）。
  \item 映射 $p$ 可能越出球外，取 $g = \frac{p}{1
        + \varepsilon/2}$。证明 $g(\bar B) \subseteq \bar B$
        且 $\sup_{\bar B}\norm{g - f} < \varepsilon$。
  \item 由第三部分导出矛盾，从而得出：
        \emph{每个连续映射 $\bar B^n \to \bar B^n$
        都有不动点。}
  \item 举例说明该定理在下列集合上不成立：开
        球；球面 $S$；闭圆环。
        各反例分别失去了 $\bar B$ 的哪一条性质？
\end{enumerate}

\medskip
\noindent\textbf{第五部分 --- 附加题。}
\begin{enumerate}[resume]
  \item （佩龙–弗罗贝尼乌斯，存在性）设 $A$ 为
        $n\times n$ 矩阵，其所有元素均 $> 0$，并设
        $\Delta = \{x \in \R^n : x_i \geq 0,\ \sum x_i =
        1\}$。证明映射 $x \mapsto Ax/\norm{Ax}_1$
        在 $\Delta$ 上良定义且连续，并且
        $\Delta$ 同胚于
        $\R^{n-1}$ 中的一个闭球（径向同胚，取自凸紧集，
        其在仿射包中具非空内部），
        从而 $A$ 有一个特征向量，
        其分量严格为正，对应特征值 $> 0$。
  \item 由此推出：每个元素均为正数的随机矩阵
        （各列之和为 $1$）都有一个平稳的
        概率向量 $\pi = A\pi$——即
        PageRank 型向量。（唯一性同样成立，但
        需要其他工具。）
  \item （毛球定理，预备）设 $v$ 在
        $S = S^{n-1}$ 的某邻域上光滑，且满足 $\langle v(x),
        x\rangle = 0$，$\norm{v(x)} = 1$，对 $x \in S$ 成立
        （称为\emph{单位切向场}）。对 $t \in \R$，令
        $F_t(x) = x + t\,v(x)$。证明
        $\norm{F_t(x)} = \sqrt{1 + t^2}$ 在 $S$ 上，$F_t$ 把
        $S$ 映入球面 $\sqrt{1+t^2}\,S$。
  \item 证明 $P(t) = \int_SF_t^*\sigma$ 是
        $t$ 的\emph{多项式}\emph{（每个系数
        函数 $F_t^*\sigma$ 在坐标卡中都是
        $t$ 的多项式，其系数关于坐标卡变量光滑；
        积分是线性的）}。
  \item 证明当 $\abs t$ 充分小时，$F_t$ 是
        从 $S$ 到 $\sqrt{1+t^2}\,S$ 上的微分同胚：
        当 $t\operatorname{Lip}(v) < 1$ 时单射；局部
        微分同胚由反函数定理得出
        （\cref{thm:b3:submanifolds:ift}，在坐标卡中）；像
        在连通的目标球面中既开又闭。
        利用\cref{lem:b3:forms:consistency}，
        以及伸缩关系 $\sigma_{\lambda x} =
        \lambda^{n}\,\sigma_x$（在 $x \mapsto \lambda x$ 下，
        需验证），得出
        \[
        P(t) = \pm(1 + t^2)^{n/2}\int_S\sigma,
        \qquad\text{符号取 } + \text{ 对较小的 } t
        \]
        （定向由 $t = 0$ 出发的连续性保持）。
  \item 由此得出（米尔诺）：若 $n$ 为\emph{奇数}，则 $(1 +
        t^2)^{n/2}$ 不是 $t$ 的多项式，但它
        与多项式 $P(t)/\int_S\sigma$ 在
        $0$ 附近相等，这是矛盾。因此\emph{
        偶数维球面 $S^{n-1}$（$n$ 为奇数）
        不带单位切向场}，并通过规范化与光滑化
        （逐分量作卷积后投影——
        两步都需说明理由），则根本不存在处处非零的
        连续切向场：地球上每股风都会留下一个
        无风点。

  \item （奇数维球面可梳理）在 $S^{2m-1}
        \subseteq \R^{2m} \cong \C^m$ 上给出一个显式光滑的
        单位切向场：在复记号下 $v(x) = \iu x$，
        即
        \[
        v(x_1, y_1, \dots, x_m, y_m) =
        (-y_1, x_1, \dots, -y_m, x_m) .
        \]
        验证其切向且为单位长，并由此得出
        第 23 题的奇偶二分法是最佳的：一个球面
        可梳理当且仅当其维数为
        奇数。第
        22 题中的多项式论证对偶数 $n$ 在何处失效？
  \item （由边界行为得零点）设 $f \colon \bar
        B^n \to \R^n$ 连续，并满足 $\langle f(x),
        x\rangle \geq 0$（对每个 $x \in S^{n-1}$）。证明
        $f$ 在 $\bar B^n$ 中某处取零值。
        \emph{（若不然，$g(x) = -f(x)/\norm{f(x)}$ 将
        $\bar B$ 连续映入 $S \subseteq \bar B$；
        对 $g$ 应用布劳威尔定理，与边界假设
        矛盾。）}由此推出\emph{满射性
        判别法}：若连续映射 $F\colon\R^n\to\R^n$ 满足
        $\frac{\langle F(x), x\rangle}{\norm x} \to
        +\infty$，当 $\norm x \to \infty$ 时，该映射是满射——
        这是非线性分析中强制性论证的
        有限维前身。
\end{enumerate}
\end{problem}
